% !TEX root = hw1-answers.tex
\chead{\textbf{Exercises week 1}}

\begin{document}

\flushleft\textbf{Topic: basic causality, taking expectations}


\section{Causal stories}
The following fictional news headlines state a correlation and suggest a causal relationship. However, we know that correlations do not always imply the obvious causation. For each headline, find three different plausible hypotheses about the causal relationships, potentially including variables that are not mentioned in the headlines. Describe the hypotheses briefly in words and with a graphical representation, such as the ones given in ``Explanation 1.1'' in the lecture notes. ($3 \times 3 \times 0.5$ pts)

\begin{enumerate}
    \item A vaccine is tested in country A and country B. In country A, patients received half a dose of the vaccine and it immunized $90\%$ of patients, while in country B, patients received the full dose and it immunized $60\%$ of patients.
    \item Of a population, 10\% belong to group A and 90\% to group B, but 90\% of those arrested for criminal offences belong to group A.
    \item High school graduation grades are highly predictive for future career success.
\end{enumerate}
\begin{gradedsolution}
Other correct answers exist. These are just examples.
\begin{enumerate}
    \item (1) Half a dose is more effective: dose $\to$ immunization. (2) the populations differed (e.g. younger in country A): country $\to$ age $\to$ protection. (3) the virus in B is more resistant than the virus in country A: country $\to$ virus $\to$ protection.
    \item (1) Group A is more predisposed to commit crime: group $\to$ predisposition to crime $\to$ crime $\to$ arrests. (2) Group A is more heavily policed so that crimes more often lead to arrests: group $\to$ policing $\to$ arrests. (3) Group A is poorer and poverty causes crime: group $\to$ poverty $\to$ crime $\to$ arrests. (4) Group A are those with a prior conviction: group $\leftarrow$ predisposition to crime $\to$ crime $\to$ arrests.
    \item (1) Working hard on high school prepares the students well for career success: HS grades $\leftarrow$ worked hard at HS $\rightarrow$ good career. (2) Students with well-educated/rich parents do well in high school and also in their career: HS grades $\leftarrow$ parents $\to$ good career. (3) employers hire from certain schools that only give high grades: HS grades $\leftarrow$ school $\to$ good career.
\end{enumerate}
\end{gradedsolution}

\section{Causal and non-causal predictions}
%Given some probabilistic graphical model, one can do probabilistic inference - using Bayes rule, marginalization and conditioning to obtain conditional probability distributions. However, some queries are causal in nature and not answerable by just probabilistic inference.
In statistics, machine learning and data science, one is often interested in making predictions based on data.
Some predictions are of a non-causal nature (for example: predict what the weather will be tomorrow),
other predictions are of a causal nature (for example: predict the effect of CO$_2$ emission on global temperature).
State three systems and for each system provide an example of a causal prediction and an example of a non-causal prediction that one could try to make based on data from the system ($3 \times 2 \times 0.5$ pts).

\begin{gradedsolution}
Not applicable
\end{gradedsolution}

\section{Discrete and Continuous Dice}
A discrete die has as sample space the six sides $\Xcal=\{a, b, c, d, e, f\}$ with uniform measure $P_D$ with probability mass function $p_D$. The value is a random variable $V_D:\Xcal \to \R$ with:
\[
    V(a)=1, V(b)=2, V(c)=3, V(d)=4, V(e)=5, V(f)=6.
\]
A continuous die has as sample space the unit interval $[0, 1]$ with uniform measure $P_C$ with probability density function $p_C$. The value is a random variable $V_C:[0,1] \to \R : x \mapsto 5 x + 1$. Using the measure theoretical notation introduced in Lecture 1, compute the expected value of the discrete and continuous dice (2+2pts).

\begin{gradedsolution}
\begin{align*}
    \E[V_D(X)]&=\int V_D dP_D= \sum_{x \in \Xcal}V_D(x)p_D(x)=6 \times 7/2 \times 1/6 = 7/2\\
    \E[V_C(X)]&=\int V_C dP_C= \int_{[0,1]}V_C(x)p(x)dx=\int_{[0,1]}V_C(x)dx \\
    &=\int_{[0,1]}5x+1dx=1+5 \times \int_{[0,1]}x dx = 1 + 5 \times 1/2=7/2
\end{align*}

\end{gradedsolution}


% POTENTIALLY USABLE FOR A LATER EXERCISE
% \section{Graphical Models \& Collider}
% It is known that men that are overweight are at high risk of contracting COVID19 and then needing intensive care. We assume that the population consists of 50\% men and 50\% woman and that half of each group is overweight. Clearly, gender is independent from being overweight. Still, when visiting an intensive care unit, one may find a higher proportion of men among the patients with a healthy weight than among the overweight patients. This is an example of the ``explain away'' effect, which occurs when the graphical model is a collider.

% We have three Boolean variables: whether the patient is male ($M$), overweight ($O$) and in need of intensive care ($C$). Given the assumption that $M$ and $O$ are independent, we can system with following graphical model:
% \begin{equation*}
%     \begin{tikzpicture}[scale=1, transform shape]
%         \tikzstyle{every node} = [draw,shape=circle]
%         \node (M) at (0,0) {$M$};
%         \node (O) at (3,0) {$O$};
%         \node (C) at (1.5,-1) {$C$};
%         \draw[-{Latex[length=3mm, width=2mm]}] (M) to (C);
%         \draw[-{Latex[length=3mm, width=2mm]}] (O) to (C);
%     \end{tikzpicture}
% \end{equation*}
% This model simply represents the fact that the joint probability $P(M, O, C)$ distribution can be factorised as:
% \[
%     P(M, O, C) = P(M)P(O)P(C|M, O)
% \]
% Because this is a collider, we have that while $M$ and $O$ are independent, $M$ and $O$ become dependent when conditioning on $C$. We study this in the particular case with conditional probability distributions:
% \begin{align*}
%     P(M=1)=P(O=1)&=1/2\\
%     P(C=1|M=0, O=0) &= 1/4\\
%     P(C=1|M=1, O=0)=P(C|M=0, O=1) &= 1/2\\
%     P(C=1|M=1, O=1)&=3/4
% \end{align*}

% Use basic probability rules to give the probability $P(M=1 | O=1, C=1)$ that an overweight patient in need of intensive care is male and the probability $P(M=1 | O=0, C=1)$ that a not-overweight patient in need of intensive care is male. (4pts)

% \begin{gradedsolution}
% \begin{align*}
% P(M=1, O=1, C=1)=3/4 \times 1/4&=3/16\\
% P(M=0, O=1, C=1)=1/2 \times 1/4&=2/16\\
% P(M=1|O=1,C=1)&=3/5\\
% P(M=1, O=0, C=1)=2/4 \times 1/4&=2/16\\
% P(M=0, O=0, C=1)=1/4 \times 1/4&=1/16\\
% P(M=1|O=0,C=1)&=2/3\\
% \end{align*}

% \end{gradedsolution}

%\end{document}
